% !CLMSynth-doc = manual
% !CLMSynth-version = 0.7.0b1
% These two lines are a machine-readable marker, so the
% file can be renamed freely, and checks the version above against the
% package version so documentation cannot silently fall behind a release.

\documentclass[a4paper,12pt]{article}
\usepackage{graphicx}
\usepackage{hyperref}
\usepackage{amsmath}
\usepackage{appendix}
\usepackage{vhistory}
\usepackage[utf8]{inputenc}
\usepackage{pdfpages}
\usepackage{listings}
\usepackage{xcolor}

\lstdefinelanguage{yaml}{
  morekeywords={true,false,null,y,n},
  sensitive=false,
  morecomment=[l]{\#},
  morestring=[b]"
}
\usepackage{booktabs}
\usepackage{tabularx}
\usepackage{amssymb}

\definecolor{codegreen}{rgb}{0,0.6,0}
\definecolor{codegray}{rgb}{0.5,0.5,0.5}
\definecolor{codepurple}{rgb}{0.58,0,0.82}
\definecolor{backcolour}{rgb}{0.95,0.95,0.92}

\lstdefinestyle{mystyle}{
    backgroundcolor=\color{backcolour},   
    commentstyle=\color{codegreen},
    keywordstyle=\color{magenta},
    numberstyle=\tiny\color{codegray},
    stringstyle=\color{codepurple},
    basicstyle=\ttfamily\footnotesize,
    breakatwhitespace=false,         
    breaklines=true,                 
    captionpos=b,                    
    keepspaces=true,                 
    numbers=none,                    
    showspaces=false,                
    showstringspaces=false,
    showtabs=false,                  
    tabsize=2
}
\lstset{style=mystyle}

\title{CLMSynth v0.7.0}
\author{Arootin Gharibian}
\date{\today}

\begin{document}

\maketitle

\begin{abstract}
CLMSynth generates a label $L(x)$ and attaches it to an existing set of cluster IDs $c(x)$ (ground truth or algorithm output), with the relationship between $L$ and $c$ specified by a configuration file and then measured afterwards. A synthetic dataset generated with this program can be used to benchmark and evaluate cluster and label matching under user-definable parameters: clusters with random class assignments and balanced or imbalanced clusters with varying degrees of cluster-label matching. Unlike existing clustering benchmarks that implicitly equate clusters with classes, the generator treats cluster formation and class assignment as independent processes with configurable matching outcomes. Each dataset instance includes ground-truth cluster membership, class labels, and metrics to create reproducible comparative benchmarking of cluster-label matching. A flexible target-metric solver can iterate over possible configurations, respecting user-defined criteria for label assignment, and solve for a CLM function within tolerance of a specific target.

\end{abstract}

\subsection*{License}
CLMSynth Documentation is licensed under \href{https://creativecommons.org/licenses/by-sa/4.0/}{Creative Commons Attribution-ShareAlike 4.0 International} license. The source code for the program is published under the MIT License at \href{https://github.com/A-Gharibian/CLMSynth}{GitHub}.

\newpage
\section{Introduction}
Existing synthetic datasets for clustering performance measurement assume equivalence between latent cluster structure and assigned class labels. The CLM dataset synthesizer addresses the natural occurrence of classes within latent clusters by treating cluster generation and class assignment as independent, configurable, and target-matched for a specific cluster-label matching. This design supports multiple benchmarks, including non-clusterable data with strong class separability, well-defined clusters with randomly assigned classes, and both balanced and imbalanced clusters. For each configuration, the output is a dataset of labels attached to the cluster membership, along with summary diagnostics characterizing cluster-label matching (CLM) measures. By generating and systematically comparing reproducible CLM scenarios, this resource complements existing clustering benchmarks and supports the evaluation of methods that study the relationships between the clustering structure and class labels.

\subsection{Scientific Objectives}
Cluster-label matching (CLM) is measured as the feature subset's effectiveness in predicting ground-truth class labels, the degree to which a feature subset supports known classes or labels regardless of clustering performance. The performance and quality outcomes associated with a selected subset of features, evaluated in terms of its ability to detect structures (clusters) and its alignment with class labels, are independent. Different algorithms and measures can optimize the selection of a given subset of features and assess how it performs in clustering, classification, and cluster-label matching. The objective of this program is to create datasets that give a user granular control over labels and match specific cluster-label matching targets, independent of the clustering, so that an algorithm or measure can reproduce specific CLM targets.

\subsection{Implementation}
The implementation in Python, together with the pipeline scripts, configuration templates, and working examples, is maintained in the \href{https://github.com/A-Gharibian/CLMSynth}{project repository}. The source code is the authoritative description of the implementation; this manual is intended for software users and includes the definition of the configuration schema and the troubleshooting reference.
\newpage

\section{Implementation Details}
\subsection{Label generator}
\begin{enumerate}
    \item \textbf{Validate}, structural checks fail fast: \texttt{num\_classes} and the data's own cluster count $K$ must each lie in $[1,64]$, a coarse backstop rather than a reliability boundary); \texttt{perfect} requires $M = K$; \texttt{single} requires \texttt{single\_match}, $M \ge 2$, $K \ge 2$; \texttt{target\_metric} is only compatible with \texttt{single}/\texttt{custom}, and \texttt{scope: pair} additionally requires \texttt{type: mcc} and \texttt{matching\_mode: single}; any requested spatial placement requires per-point coordinates; unknown modes raise errors.
    \item \textbf{Resolve proportions} $p$ are uniform if \texttt{balanced}; otherwise, explicit \texttt{proportions} take precedence, with \texttt{skew\_rule} as a fallback. Converted to integer counts by the largest-remainder rounding.
    \item \textbf{Matching mode} $\to$ alignment rules:
    \begin{itemize}
        \item \texttt{random}, a random permutation of the fixed multiset holding each label $\ell_i$ exactly $m_i$ times, rather than an independent draw from $p$ per point: the realized counts equal $m_i$, and the permutation consults neither $c(x)$ nor position, so the independence from $c(x)$ holds by construction rather than expectation.
        \item \texttt{perfect}, bijection $\pi: \text{cluster} \to \text{label}$ in sorted order; label counts are assigned to the paired cluster sizes; $L(x) = \pi(c(x))$.
        \item \texttt{single}, one target pair $(k^*, \ell^*)$ aligned at full recall; all other points are assigned by spillover.
        \item \texttt{custom}, an \texttt{assignment\_matrix} of rules, each a \texttt{recall\_target} fraction of one label's point budget into a chosen set of clusters (enabling surjective, partial, and overlapping).
    \end{itemize}
    \item \textbf{Allocation \& feasibility}, each rule is divided into its clusters (\texttt{split\_rule}); exceeding the cluster capacity raises an infeasibility error stating the maximum feasible recall. The remaining cluster capacity is filled by \texttt{spillover\_rule}. (\texttt{proportional\_to\_marginal} preserves the target proportions).
    
    \item \textbf{Target metric}: the target metric and the solver vary only the recall level; every setting above is held fixed across probes, so the noise structure changes the solved recall rather than being applied after it.

With \texttt{scope: global} (the default), one recall level is solved for the whole partition by a coarse grid scan; with \texttt{scope: pair}, the single-pair MCC of the \texttt{single\_match} pair is inverted in closed form and no search runs. A target above the configuration's reachable ceiling terminates at the closest achievable value, and a target below its floor terminates against that in the same way. Because $R_K$ is a maximum over matching, and the mean of $R_K$ over all relabeling is exactly $0$, the matched value is never negative and its expectation under independent labeling is strictly positive. Therefore, a requested $0.0$ is not null.

The engine measures CLM, and a solution that converged internally can still land outside the tolerance, most visible at small $N$, where the placement of spillovers is a large part of the result.
    
Currently, three CLM measures are available: \texttt{clustering\_mcc} is the permutation-invariant \texttt{scikit-learn} implementation of Gorodkin's  $R_K$, which is a one-to-one matching of the cluster--label contingency table that maximizes $R_K$ as the value solved for.\\
\texttt{scope: global}. \texttt{clustering\_mcc\_pair} is the $2\times2$ Matthews $\phi$ of a cluster against a label, which does not need matching.\\
\texttt{clustering\_ari} is solved for the adjusted Rand inde (ARI): The ARI measures the agreement over \emph{pairs} of points rather than through any correspondence between the label and cluster identifiers, and is chance-corrected so that independent partitions score $0$. Because it is a function of the cell of the contingency table and the marginal pair counts alone, it does not need an alignment step: permuting the values of the generated label leaves it unchanged, it is symmetric in its two arguments, and it is defined for any $M$ and $K$. This is what distinguishes it from $R_K$, which is meaningful only after the Hungarian alignment.
    
    \item \textbf{Spatial placement} Configures optionally, within each cluster, \emph{which} points receive each label is chosen by weighted sampling without replacement, with weights derived from normalized distance to the cluster centroid, with\texttt{profile}: \texttt{linear}/\texttt{step}/\texttt{exponential}; and \texttt{favors}: \texttt{steepness} for exponential or \texttt{core}/\texttt{boundary} for other profiles. The contingency table, and therefore the target CLM outcome, is fixed by stages 2--4, so spatial structure can be varied at a constant metric value.

\end{enumerate}

\noindent\textbf{Note on multi-class multi-label measurement:} Releases of up to 0.6.9 were matched on the grid-search to maximize precision, and thus reported the $R_K$ of accuracy-optimal matching. The two agree where the rules impose real alignment and differ only near the lower end of the target metric, specifically for unbalanced datasets.
\newpage

\subsection{Configuration file definition}
\noindent The \texttt{clm\_label\_engine} is the core module for label synthesis. The configuration parameters for the core module are defined as follows:
\begin{lstlisting}[language=yaml]
clm_label:
  num_classes: M                   # 1 <= M <= 64; the data's own K is capped the
                                   # same way (CLM-126 / CLM-127)
  balance: balanced | unbalanced   # balanced: uniform 1/M, proportions ignored
  proportions: [...]               # unbalanced: used directly when given. Must have
                                   # exactly M entries (CLM-121) summing to 1 (CLM-106)
  skew_rule: geometric | dominant_minority | dirichlet
                                   # unbalanced fallback when no proportions
  skew_params:
    ratio: 0.5                     # geometric: p_i ~ ratio^i
    dominant_index: 0              # dominant_minority
    dominant_share: 0.5            # dominant_minority
    alpha: 1.0                     # dirichlet: Dirichlet(alpha,...,alpha)
  matching_mode: perfect | single | random | custom
  single_match: {cluster: k*, label: l*}    # required if matching_mode: single
  assignment_matrix:                         # required if matching_mode: custom
    - {label: l, clusters: [k1, ...], recall_target: r}
  split_rule: proportional_to_size | equal   # how a rule's mass splits over its clusters
  spillover_rule: proportional_to_marginal | uniform | concentrated
  concentrated_labels: [l, ...]    # spillover_rule: concentrated only. A LIST of
                                   # integer ids in 0..M-1 (CLM-128); a bare number
                                   # is read as a range, not a one-element list
                                   # (default: the single largest label)
  target_metric:                   # optional (single/custom only)
    type: mcc | ari
    value: 0.7
    scope: global | pair           # global (default): one recall solved
                                   #   numerically for the whole partition.
                                   # pair: exact closed form for the
                                   #   single_match pair; requires type: mcc
                                   #   and matching_mode: single.
    tolerance: 0.01                # both scopes: how far the DELIVERED
                                   #   labeling may sit from the target
                                   #   before CLM-309/CLM-310 is raised
    max_iter: 40                   # global only, bisection budget; the pair
                                   #   scope inverts exactly and never iterates
    probe_seed: 42                 # optional; defaults to the run seed
  competing_noise:                 # optional (single/custom): structured noise.
                                   # Converts a share of ONE cluster's UNCLAIMED
                                   # points to one competing label. Bypasses the
                                   # target proportions by design (CLM-304).
    - cluster: k
      label: l_comp
      share: 1.0                   # fraction of that cluster's leftover points
      favors: boundary | core | random
  centroid_dependence:
    enabled: true | false
    profile: linear | exponential | step
    favors: core | boundary
    steepness: 3.0                 # exponential only
\end{lstlisting}

\noindent Spatial placement (\texttt{centroid\_dependence} or a \texttt{competing} entry \_noise \texttt{that favors the core}/\texttt{boundary}) requires vectors of characteristic per-point. The complete catalog of coded errors and warnings is the reference for troubleshooting.

\newpage
\section{Usage}
The input variables required for the generator (\texttt{X} and \texttt{c}; cluster centroids are derived internally from \texttt{X}) are expected to be supplied by external clustered datasets. CLMSynth is designed to work with online and offline datasets, such as the Gagolewski's clustering benchmark suite (\texttt{clustbench}) or synthetic datasets generated via MDCGen (\texttt{mdcgen}). In addition, synthetic data generated through the offline generator (\texttt{fabricated\_data}) and direct import through \texttt{byoc} (\emph{bring-your-own-clusters}) are possible.

\noindent The package installs three entry points: \texttt{clmsynth} runs the pipeline with a YAML configuration, \texttt{clmsynth-config} renders a configuration from an upstream payload, and \texttt{clmsynth-wizard} generates one interactively. A simple demonstrative run is as follows:

\begin{lstlisting}[language=Python]
# Install package with requirements, then:
from clmsynth.clm_label_engine import generate_clm_labels
# An arbitrary clustering data:
rng = np.random.default_rng(0)
sizes = [200, 150, 150]
c = np.repeat([0, 1, 2], sizes)
X = np.vstack([rng.normal(mu, 1.0, size=(n, 2))
               for mu, n in zip([(0, 0), (6, 0), (0, 6)], sizes)])
# Basic configuration:
config = {
    "num_classes": 3,
    "balance": "unbalanced",
    "proportions": [0.5, 0.3, 0.2],
    "matching_mode": "custom",
    "assignment_matrix": [
        {"clusters": [1, 2], "label": 0, "recall_target": 0.9},
    ],
    "split_rule": "proportional_to_size",
    "spillover_rule": "proportional_to_marginal",
    "centroid_dependence": {
        "enabled": True,
        "profile": "exponential",
        "favors": "core",
        "steepness": 4.0,
    },
}
# Generate labels:
L = generate_clm_labels(cluster_labels=c, coords=X, cfg=config, seed=42)
print(L.value_counts().sort_index().to_dict())
# -> {0: 250, 1: 150, 2: 100}
\end{lstlisting}

\newpage

\section{Example Outputs}
The global $R_K$/ARI between an $M$-label and a $K$-cluster partition has no closed form, which is why \texttt{scope: global} is solved numerically. The compact closed-form expression $TP^\star$ inverts the \emph{single} $2\times2$ MCC (what \texttt{scope: pair} evaluates), which is why that mode requires \texttt{type: mcc} and \texttt{matching\_mode: single} and reaches its target without a search (see the published article). That is why in table~\ref{tab:g2mg-outcomes} The CLM target of single-match configurations A3 and D1 is feasible. Configurations A4 and B1--B3 share one set of rules (labels sized $[1024,614,410]$, one rule per label); B1 and B2 request possible values within their configuration (solver within the tolerance $0.01$, while B3, a 3-label partition of a 2-cluster, requests $\mathrm{MCC}=0.99$, which is not possible. The solver returns its closest feasible value, $0.718$, which is the maximum attainable for this configuration. C1--C3 differ only in spatial placement.

\begin{table}[!h]
\centering\tiny
\begin{tabular}{@{}l p{3.5cm} p{4.0cm} p{4.6cm}@{}}
\toprule
\textbf{Cfg} & \textbf{Configuration} & \textbf{Expected} & \textbf{Actual} \\
\midrule
A1 & perfect, $M{=}K{=}2$ & bijection; MCC $\approx 1$ & MCC $=1.000$, ARI $=1.000$\\
A2 & random, $M{=}4$ & no cluster info; ARI $\approx 0$, $R_K$ at its null level & MCC $=0.018$, ARI $=-0.000$\\
A3 & single (cluster 1, $M{=}4$) & feasible here ($512<1024$); partial alignment & MCC $=0.279$, ARI $=0.166$\\
A4 & custom, explicit recalls (0.8/0.5/0.3) & feasible; counts $[1024,614,410]$ & MCC $=0.527$, ARI $=0.393$\\
B1 & custom + solve MCC$=0.6$ & reachable here; solver hits it within tol. & solved MCC $=0.598$ (target 0.6)\\
B2 & custom + solve ARI$=0.4$ & reachable here; solver hits it within tol. & solved ARI $=0.391$ (target 0.4)\\
B3 & custom + solve MCC$=0.99$, max\_iter 15 & unreachable; non-convergence + delivered-value check & solved MCC $=0.718$\\
C1 & single largest + centroid linear/core & MCC $=$ A3 (placement in 4-D) & MCC $=0.279$, ARI $=0.166$\\
C2 & single largest + centroid exp/boundary & same contingency as C1 & MCC $=0.279$, ARI $=0.166$\\
C3 & single largest + centroid step/core & same contingency as C1 & MCC $=0.279$, ARI $=0.166$\\
D1 & single (cluster 1, $M{=}2$) & fills cluster exactly; MCC $\approx 1$ & MCC $=1.000$, ARI $=1.000$\\
\bottomrule
\end{tabular}
\caption{Cluster--label matching outcomes on \texttt{g2mg\_4\_20} ($N=2048$, 4-D, $K=2$ clusters of 1024 points, reference labeling \texttt{labels0}).}
\label{tab:g2mg-outcomes}
\end{table}

\newpage

\begin{versionhistory}
  \vhEntry{0.1}{2026-07-07}{AG}{Initial release: CLM label engine, \texttt{[CLM-\#\#\#]} diagnostics, configuration wizard, MCC/ARI evaluation.}
  \vhEntry{0.2}{2026-07-15}{AG}{Internal, unpublished.}
  \vhEntry{0.3}{2026-07-31}{AG}{\texttt{clustering\_mcc\_pair} added to the API; \texttt{fabricated\_data} cluster ids changed from strings to integers.}
  \vhEntry{0.4}{2026-07-31}{AG}{The renderer can now emit \texttt{target\_metric}, \texttt{competing\_noise}, \texttt{concentrated\_labels}, \texttt{skew\_params} and \texttt{steepness}; \texttt{mdcgen} source repaired.}
  \vhEntry{0.5}{2026-07-31}{AG}{Research release: \texttt{CITATION.cff}, \texttt{[CLM-1xx]} configuration errors abort the run.}
  \vhEntry{0.6.2}{2026-08-02}{AG}{First public release. Label-space and pair-scope guards: \texttt{[CLM-121]}, \texttt{[CLM-128]}, \texttt{[CLM-129]}, \texttt{[CLM-130]}, \texttt{[CLM-153]} and \texttt{[CLM-310]} reject configurations that previously produced silently incorrect labels; documentation revised throughout}
  \vhEntry{0.6.3}{2026-08-08}{AG}{Correctness release: \texttt{[CLM-131]} validates \texttt{skew\_params}; \texttt{scope: pair} now honors \texttt{tolerance}; run folders are created atomically; \texttt{[CLM-104]}/\texttt{[CLM-105]} report per dataset instead of aborting a batch.}
  \vhEntry{0.6.4}{2026-08-09}{AG}{Tooling and import hygiene: the target-metric search and the final generation now share one random stream, so a solved value is the value delivered; \texttt{byoc} imports are validated against stated requirements (see the uncoded-rejections section of the troubleshooting reference.}
  \vhEntry{0.6.5}{2026-08-10}{AG}{The test suite becomes part of the repository.}
  \vhEntry{0.6.6}{2026-08-11}{AG}{A patch release}
  \vhEntry{0.6.7}{2026-08-12}{AG}{CLI-Wizard update and improvement}
  \vhEntry{0.6.8}{2026-08-20}{AG}{Patch without release, preparation for CLMSynth-GUI}
  \vhEntry{0.6.9}{2026-08-23}{AG}{Patch release, staging for release 0.7.0}
  \vhEntry{0.7.0b1}{2026-08-30}{AG}{Article-review release: maximize the multiclass \texttt{R\_K} rather than accuracy; Brent's method for the recall solver ; \texttt{byoc\_suite.tag\_columns} carries non-feature columns through an import; support for Python 3.11 dropped.}
\end{versionhistory}
\newpage
% \section*{License}

% \begin{verbatim}
% The MIT License

% Copyright (c) 2026 Arootin Gharibian <arootin.gharibian@vsb.cz>
% \end{verbatim}
% Permission is hereby granted, free of charge, to any person obtaining a copy of this Software and associated documentation files (the "Software"), to deal in the Software without restriction, including without limitation the rights
% to use, copy, modify, merge, publish, distribute, sublicense, and/or sell copies of the Software, and to permit persons to whom the Software is furnished to do so, subject to the following conditions:

% The above copyright notice and this permission notice shall be included in all copies or substantial portions of the Software.

% THE SOFTWARE IS PROVIDED "AS IS", WITHOUT WARRANTY OF ANY KIND, EXPRESS OR IMPLIED, INCLUDING BUT NOT LIMITED TO THE WARRANTIES OF MERCHANTABILITY, FITNESS FOR A PARTICULAR PURPOSE AND NONINFRINGEMENT. IN NO EVENT SHALL THE
% AUTHORS OR COPYRIGHT HOLDERS BE LIABLE FOR ANY CLAIM, DAMAGES OR OTHER LIABILITY, WHETHER IN AN ACTION OF CONTRACT, TORT OR OTHERWISE, ARISING FROM, OUT OF OR IN CONNECTION WITH THE SOFTWARE OR THE USE OR OTHER DEALINGS IN THE SOFTWARE.


\section*{Licenses for Third-Party Libraries}

Python is distributed under the Python Software Foundation License V.2. This project relies on several open-source third-party Python libraries, including \texttt{NumPy}, \texttt{Pandas}, \texttt{SciPy}, \texttt{scikit-learn}, \texttt{Matplotlib}, \texttt{Seaborn}, and \texttt{PyYAML}. Optionally \texttt{faker} and \texttt{mdcgenpy}. These packages are distributed under their respective licenses (BSD 3-Clause, MIT, and Apache 2.0). Refer to the official documentation of each package or the project repository for full third-party license details.

\section{References}
\begin{itemize}
    \item Gagolewski, M. (2020). \textit{A framework for benchmarking clustering algorithms}. \url{https://clustering-benchmarks.gagolewski.com/weave/suite-v1.html}
    \item Iglesias, F., Zseby, T., Ferreira, D. \textit{et al.} (2019). \textit{MDCGen: A Multidimensional Dataset Generator for Clustering}. Journal of Classification. \url{https://link.springer.com/article/10.1007/s00357-019-9312-3}
    \item Jeon, H., Aupetit, M., Shin, D., Cho, A., Park, S., and Seo, J. (2025). \textit{Measuring the Validity of Clustering Validation Datasets}. IEEE Transactions on Pattern Analysis and Machine Intelligence, 47(6), 5045-5058. \url{https://doi.org/10.1109/TPAMI.2025.3548011}
\end{itemize}



\includepdf[pages=-]{CLMSynth_configuration_troubleshooting.pdf}
\includepdf[pages=-]{engine_decision_map_v068.pdf}
\end{document}
